\begin{abstract}

We propose the \DH{} as a structural framework for reinterpreting
divergence, renormalization, and the limits of physical description.
Rather than treating divergence only as a technical failure to be corrected,
we read it as a possible boundary signal:
an indication that a descriptive window has been pushed toward the limit
of the differences it can represent as finite quantities.

From this perspective,
renormalization is not an operation that erases a boundary.
It is a finite re-expression of boundary contact inside a chosen
observational window.
The speed-of-light limit,
the Planck-scale breakdown of spacetime continuity,
and black-hole singularities are re-positioned as local or hierarchical
horizon-forms,
rather than as unrelated failures or as different objects beyond theory.

This framework does not propose a unified theory of interactions.
Instead,
it identifies a prior structural layer:
the condition under which registered difference,
trace formation,
and physical description become possible.
The \DHP{} states that a description based on difference cannot convert
the absence of difference into an ordinary observable state.

\VED{} (Vortical Enclosure Dynamics) is positioned within this framework
as an open generative description.
It does not explain existence from outside existence.
It describes why generated spacetime encounters a descriptive boundary
when it attempts to reach the condition from which its own differences arise.

\end{abstract}
